% Options for packages loaded elsewhere
% Options for packages loaded elsewhere
\PassOptionsToPackage{unicode}{hyperref}
\PassOptionsToPackage{hyphens}{url}
%
\documentclass[
  11pt,
  letterpaper,
  DIV=11,
  numbers=noendperiod,
  twoside]{scrartcl}
\usepackage{xcolor}
\usepackage[left=1.1in, right=1in, top=0.8in, bottom=0.8in,
paperheight=9.5in, paperwidth=7in, includemp=TRUE, marginparwidth=0in,
marginparsep=0in]{geometry}
\usepackage{amsmath,amssymb}
\setcounter{secnumdepth}{3}
\usepackage{iftex}
\ifPDFTeX
  \usepackage[T1]{fontenc}
  \usepackage[utf8]{inputenc}
  \usepackage{textcomp} % provide euro and other symbols
\else % if luatex or xetex
  \usepackage{unicode-math} % this also loads fontspec
  \defaultfontfeatures{Scale=MatchLowercase}
  \defaultfontfeatures[\rmfamily]{Ligatures=TeX,Scale=1}
\fi
\usepackage{lmodern}
\ifPDFTeX\else
  % xetex/luatex font selection
  \setmathfont[]{Garamond-Math}
\fi
% Use upquote if available, for straight quotes in verbatim environments
\IfFileExists{upquote.sty}{\usepackage{upquote}}{}
\IfFileExists{microtype.sty}{% use microtype if available
  \usepackage[]{microtype}
  \UseMicrotypeSet[protrusion]{basicmath} % disable protrusion for tt fonts
}{}
\usepackage{setspace}
% Make \paragraph and \subparagraph free-standing
\makeatletter
\ifx\paragraph\undefined\else
  \let\oldparagraph\paragraph
  \renewcommand{\paragraph}{
    \@ifstar
      \xxxParagraphStar
      \xxxParagraphNoStar
  }
  \newcommand{\xxxParagraphStar}[1]{\oldparagraph*{#1}\mbox{}}
  \newcommand{\xxxParagraphNoStar}[1]{\oldparagraph{#1}\mbox{}}
\fi
\ifx\subparagraph\undefined\else
  \let\oldsubparagraph\subparagraph
  \renewcommand{\subparagraph}{
    \@ifstar
      \xxxSubParagraphStar
      \xxxSubParagraphNoStar
  }
  \newcommand{\xxxSubParagraphStar}[1]{\oldsubparagraph*{#1}\mbox{}}
  \newcommand{\xxxSubParagraphNoStar}[1]{\oldsubparagraph{#1}\mbox{}}
\fi
\makeatother


\usepackage{longtable,booktabs,array}
\newcounter{none} % for unnumbered tables
\usepackage{calc} % for calculating minipage widths
% Correct order of tables after \paragraph or \subparagraph
\usepackage{etoolbox}
\makeatletter
\patchcmd\longtable{\par}{\if@noskipsec\mbox{}\fi\par}{}{}
\makeatother
% Allow footnotes in longtable head/foot
\IfFileExists{footnotehyper.sty}{\usepackage{footnotehyper}}{\usepackage{footnote}}
\makesavenoteenv{longtable}
\usepackage{graphicx}
\makeatletter
\newsavebox\pandoc@box
\newcommand*\pandocbounded[1]{% scales image to fit in text height/width
  \sbox\pandoc@box{#1}%
  \Gscale@div\@tempa{\textheight}{\dimexpr\ht\pandoc@box+\dp\pandoc@box\relax}%
  \Gscale@div\@tempb{\linewidth}{\wd\pandoc@box}%
  \ifdim\@tempb\p@<\@tempa\p@\let\@tempa\@tempb\fi% select the smaller of both
  \ifdim\@tempa\p@<\p@\scalebox{\@tempa}{\usebox\pandoc@box}%
  \else\usebox{\pandoc@box}%
  \fi%
}
% Set default figure placement to htbp
\def\fps@figure{htbp}
\makeatother





\setlength{\emergencystretch}{3em} % prevent overfull lines

\providecommand{\tightlist}{%
  \setlength{\itemsep}{0pt}\setlength{\parskip}{0pt}}



 


\setlength\heavyrulewidth{0ex}
\setlength\lightrulewidth{0ex}
\usepackage[automark]{scrlayer-scrpage}
\clearpairofpagestyles
\cehead{
  Brian Weatherson
  }
\cohead{
  Logical Positivism and The Rise of Analytic Philosophy
  }
\ohead{\bfseries \pagemark}
\cfoot{}
\makeatletter
\newcommand*\NoIndentAfterEnv[1]{%
  \AfterEndEnvironment{#1}{\par\@afterindentfalse\@afterheading}}
\makeatother
\NoIndentAfterEnv{itemize}
\NoIndentAfterEnv{enumerate}
\NoIndentAfterEnv{description}
\NoIndentAfterEnv{quote}
\NoIndentAfterEnv{equation}
\NoIndentAfterEnv{longtable}
\NoIndentAfterEnv{abstract}
\renewenvironment{abstract}
 {\vspace{-1.25cm}
 \quotation\small\noindent\emph{Abstract}:}
 {\endquotation}
\newfontfamily\tfont{EB Garamond}
\addtokomafont{disposition}{\rmfamily}
\addtokomafont{descriptionlabel}{\rmfamily}
\addtokomafont{title}{\normalfont\itshape}
\let\footnoterule\relax

\makeatletter
\renewcommand{\@maketitle}{%
  \newpage
  \null
  \vskip 2em%
  \begin{center}%
  \let \footnote \thanks
    {\itshape\huge\@title \par}%
    \vskip 0.5em%  % Reduced from default
    {\large
      \lineskip 0.3em%  % Reduced from default 0.5em
      \begin{tabular}[t]{c}%
        \@author
      \end{tabular}\par}%
    \vskip 0.5em%  % Reduced from default
    {\large \@date}%
  \end{center}%
  \par
  }
\makeatother
\RequirePackage{lettrine}

\renewenvironment{abstract}
 {\quotation\small\noindent\emph{Abstract}:}
 {\endquotation\vspace{-0.02cm}}

\setmainfont{EB Garamond Math}[
  BoldFont = {EB Garamond SemiBold},
  ItalicFont = {EB Garamond Italic},
  RawFeature = {+smcp},
]

\newfontfamily\scfont{EB Garamond Regular}[RawFeature=+smcp]
\renewcommand{\textsc}[1]{{\scfont #1}}

\renewcommand{\LettrineTextFont}{\scfont}
\KOMAoption{captions}{tableheading}
\makeatletter
\@ifpackageloaded{caption}{}{\usepackage{caption}}
\AtBeginDocument{%
\ifdefined\contentsname
  \renewcommand*\contentsname{Table of contents}
\else
  \newcommand\contentsname{Table of contents}
\fi
\ifdefined\listfigurename
  \renewcommand*\listfigurename{List of Figures}
\else
  \newcommand\listfigurename{List of Figures}
\fi
\ifdefined\listtablename
  \renewcommand*\listtablename{List of Tables}
\else
  \newcommand\listtablename{List of Tables}
\fi
\ifdefined\figurename
  \renewcommand*\figurename{Figure}
\else
  \newcommand\figurename{Figure}
\fi
\ifdefined\tablename
  \renewcommand*\tablename{Table}
\else
  \newcommand\tablename{Table}
\fi
}
\@ifpackageloaded{float}{}{\usepackage{float}}
\floatstyle{ruled}
\@ifundefined{c@chapter}{\newfloat{codelisting}{h}{lop}}{\newfloat{codelisting}{h}{lop}[chapter]}
\floatname{codelisting}{Listing}
\newcommand*\listoflistings{\listof{codelisting}{List of Listings}}
\makeatother
\makeatletter
\makeatother
\makeatletter
\@ifpackageloaded{caption}{}{\usepackage{caption}}
\@ifpackageloaded{subcaption}{}{\usepackage{subcaption}}
\makeatother
\usepackage{bookmark}
\IfFileExists{xurl.sty}{\usepackage{xurl}}{} % add URL line breaks if available
\urlstyle{same}
% fallback for those not using the hyperref driver hyperxmp:
\makeatletter
\@ifundefined{xmpquote}{\newcommand{\xmpquote}[1]{#1}}{}
\makeatother
\hypersetup{
  pdftitle={Logical Positivism and The Rise of Analytic Philosophy},
  pdfauthor={Brian Weatherson},
  hidelinks,
  pdfcreator={LaTeX via pandoc}}


\title{Logical Positivism and The Rise of Analytic Philosophy}
\author{Brian Weatherson}
\date{2026}
\begin{document}
\maketitle
\begin{abstract}
Contribution to the book launch event on Sander Verhaegh's \emph{Logical
Positivism}.
\end{abstract}


\setstretch{1.1}
\emph{Logical Positivism} is a great book. It combines careful attention
to archival records, quantitative data, and social context, to tell an
incredibly rich story about the rise, and I guess fall, of logical
positivism in America. I'm going to recommend it to all my students as
crucial background for how philosophy in America got to where it is, and
a correction to several popular narratives.

On the data side, this is one of the best examples you'll see of data
use that highlights simple understandable statistics, the main one in
the book is simply how often philosophers are mentioned, with awareness
of the limits of these statistics. I'll just mention one aspect of this
that stuck out. The book often presents tables of outliers, e.g.,
philosophers who are mentioned most often or (relative to expectation)
least often in a particular corpus. Whenever I build tables like that,
they are mostly records of where the data needs cleaning. These tables
look simple, but the work that goes into making the outliers look
meaningful is so demanding, and the results here are fascinating.

The book mentions the social impact of positivism, this is a big focus
of chapter 11, and there's a question here I'd love to know. Maybe we
could even throw quantitative tools at it. Is there any documented line
from positivism to the (really rather terrible) fact/opinion worksheets
American children still get in elementary school? They \emph{look}
positivist to me, but it's possible that the origin comes from
pragmatism, and especially Dewey. The conservative complaints in chapter
11 that positivists were separating facts and values would look really
weird to students educated in the US over the last couple of
generations.

I found especially pleasing in the book was the role that \emph{public}
universities played in the story. The focus is so often on the famous
Ivy League schools, and five of them are central to Verhaegh's book. But
the story doesn't work without the public universities too. We see
Schlick get a huge audience at Berkeley very early on, then Feigl gets
hired by Iowa and given more resources at Minnesota. Reichenbach, after
many challenges, is hired at UCLA. The central university at the end of
the narrative, and still the home of the most important positivist
archive, is Pitt. And after Charles Stevenson is denied tenure at Yale,
he gets a permanent home at Michigan, where I'm based. To this day, our
most prestigious research prize for graduate students is the Stevenson
prize.

Chapter 9 is on the rise of `analytic philosophy', and it's possible
here I slightly disagreed with Verhaegh, at least in emphasis. I came
away from this chapter thinking that the rise was \emph{somewhat}
inevitable, at least from a few years before it happened. We're taught
to distrust inevitability narratives, with good reason, and while
everything I'm about to say is in the book, I might be pushing the
narrative further than he intended.

Verhaegh's narrative about the rise of analytic philosophy starts, as it
really has to, with the tricky question of what analytic philosophy
\emph{is}. Famously, this is far from obvious. All philosophical
traditions develop varieties; with analytic philosophy that variety was
there at the start. Just what do Carnap, Moore, and Anscombe, for
example, have in common?

The story here Verhaegh tells is fascinating in two respects. First, the
label `analytic philosophy' came before there was anything for it to
clearly label. This was a stellar example of what we now call conceptual
engineering. What it came to label was the union of positivism, broadly
construed, with realism, especially Cambridge realism. This wasn't the
only label, `scientific philosophy' was also used, but it was a label,
and it ultimately stuck. \emph{After} the label was applied, this became
a movement.

This movement wasn't unified theoretically; the analytics disagreed
\emph{deeply} with other analytics. Think Austin vs Ayer on perception,
or Anscombe and Foot vs Hare on moral realism. But these analytics hired
each other, they published each other, and they built schools where one
could study both positivism and realism.

As well as these points of institutional union, they had two other
commonalities. One was about the importance of Russell. As Verhaegh
points out, this was important to the positivists. By aligning
themselves with a movement started by Russell, they acquired some
longevity; they were no longer a passing fad. That Russell had moved
away from realism towards more positivist views between the wars mustn't
have hurt. The other was that logic and philosophy of science were
important parts of philosophy. Those two commonalities turned out to be
enough.

This positivism/realism alliance was not inevitable. A simple story of
pre-war American philosophy was that there were four big schools:
idealism, pragmatism, realism, and positivism. There were others, but
those stood out. Doctrinally, of the other three schools, realism seemed
the least likely ally for positivism. There are many worlds where this
alliance doesn't get formed; if Ramsey lives it's really easy to see a
pragmatism/positivism alliance, for instance. But it did get formed. And
from that point I feel the success was basically inevitable.

One reason for that, which Verhaegh stresses, is that it couldn't help
grow. The GI Bill then the baby boom led to huge amounts of hiring. The
success of science in the war, then the space race, meant that this
would involve lots of maths and science. The faction that had what
everyone was bidding for was bound to grow.

But the other reason is that the other side is a mess, and it's really
hard to see \emph{how} it stops the rise of the analytics. If positivism
had arisen a generation earlier, it may have struggled to dislodge
idealism from its central place in Britain and America. But by 1950,
idealism had basically died. None of the other
non-analytic/anti-analytic movements was strong enough on its own to be
a serious rival. Verhaegh notes that Columbia, the home of Deweyan
pragmatism/naturalism, was basically only hiring Columbia grads. I see
that as a sign of weakness; if there were pragmatist students
everywhere, some would be so good you simply had to hire them. There
just weren't.

If the second reason is right, the way to stop it would have been for
the rival forces, the `humanist philosophies' to have allied
institutionally, the way positivists and realists did. That's how there
was, eventually, a rival to analytic philosophy in the form of
`continental philosophy'. That too is a hodge-podge of theories and
theorists, but it's a hodge-podge that worked together. Could the
`humanists' have done something similar? They did have a name, `humanist
philosophy', but in practice it never formed a school, or a super-school
like analytic and continental. That's for a good reason: the
anti-analytic forces at various schools were so diverse. Verhaegh
stresses the metaphysicians at Yale, and pragmatists at Columbia, and
Christian humanists at Princeton. If they had formed an institutional
alliance, hiring each other and building a curriculum where students
studied all of them, maybe that could have stopped the rise. But nothing
like that came close to happening.

So I'll end with a question, one that I guess relates to the chapter 13
discussion of explanation. I think counterfactuals are crucial to
explanation. If we want to explain why analytic philosophy became
dominant, we need to ask what the philosophical world would have looked
like if it never became dominant. My question is, what does that world
look like? And how late can it reasonably, feasibly, differ from actual
history? My guess is that the deviation from actual history has to come
relatively early; by 1950, the outcome is baked in, even though the
analytics don't actually take over for several more years. But I'd like
to be convinced I'm wrong; I get nervous with inevitability claims!

To wrap up, this is a fantastic book, full of insights on both the
history of philosophy, and the methodology of that history. I hope it's
really widely read.



\vspace{12pt} \noindent Draft.


\end{document}
